\begin{abstract}
The ProveIt collection contains dozens of detailed finitary reconstructions of
mathematical arguments, together with several general frameworks based on finite
observations, contracts, demand propagation, relational witnesses, and stable
observables.  The persistent gap is now sharply identifiable: a terminating
algorithm or a host-level Lean theorem is not yet an object-language derivation
\[
             \Prf\;\PRAAx\;[]\;\Phi.
\]
This report proposes a general framework that targets the repository's
kernel-checked five-rule clause calculus \(\KPrf\) as an intermediate proof
language.  A mathematical development is compiled into a finite \emph{native
clause program}: every node exports a quantifier-free native formula or clause,
every edge is an explicit finite dependency, and every node carries a local
\(\KPrf\) derivation.  Program composition is then not justified by semantic
reflection; it is ordinary \(\KPrf\) cut.  Parameter specialization is ordinary
\(\KPrf\) substitution.  The checked map \(\KPrf\to\NPrf\) followed by
\path{native_iff_pra} exports the result to ordinary iPRA.  Conversely, checked
\(\Prf\;\PRAAx\) derivations---notably the new polynomial identity backend in
\path{IPRADerivations}---can be imported as native clause nodes.

The report develops full human proofs of the compiler theorem, bounded trace
reflection, least bounded search, acyclic backward demand propagation, bounded
least fixed points for cyclic precision constraints, relational viability,
stable-observer elimination of noncanonical witnesses, exact recovery under a
positive gap, finite exceptional-set synchronization, lexicographic descent
fuel, and the equivalence between an explicit code-height bound and complete
finite enumeration.  It gives reusable presentations for integers, rationals,
represented reals and complex numbers, finite algebraic structures, matrices,
polynomials, simplicial data, and finite categories.  Representative
applications explain how the method handles polynomial identities,
fixed-dimensional Cayley--Hamilton, numerical analysis with moduli, root
clusters, spectral labels, finite topology, and conditional Diophantine
arguments.

The framework resolves several recurring repository difficulties: hidden host
computation, substitution bookkeeping, graph composition, changing precision,
cyclic resource requirements, and unstable witness selection.  It does not
manufacture absent deep quantitative producers.  General effective height
bounds, uniform Mordell--Weil covering representatives, ineffective
prime-distribution constants, and singular exact branch selectors remain
explicit interfaces rather than being renamed as certificates.
\end{abstract}

\section*{Executive conclusions}

The central recommendation is to standardize future developments on a
\emph{two-backend proof pipeline}.

\begin{enumerate}[label=\textbf{E\arabic*.}]
\item \textbf{Use ordinary \(\Prf\;\PRAAx\) for domain algebra.}
The new \path{IPRADerivations} library already supplies checked arithmetic,
semiring, integer, and polynomial-identity infrastructure.  A ring identity
should be proved there rather than reimplemented inside every packet.

\item \textbf{Use \(\KPrf\) for finite control and composition.}
The clause calculus is closed under substitution and has a proved cut theorem.
It is therefore a better intermediate representation for bounded loops,
search, finite branching, and proof graphs than an ad hoc host structure whose
soundness is still to be formalized.

\item \textbf{Compile concrete syntax, not encoded programs.}
A theorem-specific finite graph is specialized to actual terms, formulas, and
proof constructors before it is internalized.  No total primitive-recursive
universal evaluator and no truth-to-provability reflection principle is used.

\item \textbf{Keep relations when selection is unstable.}
Approximate roots, local charts, spectral subspaces, and compatible finite
models are represented by finite witness clouds.  A selector is introduced only
when bounded viability is needed; branch-independent information is extracted
by a stable observer without pretending that a canonical witness exists.

\item \textbf{Make quantitative producers visible.}
A search cutoff, convergence modulus, denominator bound, separation gap,
finite-cover certificate, or height bound is part of the theorem statement or
is produced by a separately proved node.  A cardinality statement is not a
search bound, and a terminating metaprogram is not an iPRA derivation.
\end{enumerate}

The package includes a dependency-ordered Lean plan, candidate bridge code,
exact finite reference algorithms, fourteen passing regression tests, a theorem
ledger, and a difficulty matrix.  The intended first implementation milestone
is small: kernel-check the list-cut and graph-fold layer against current
\path{NativeClausalDiagonal.lean}; then connect one existing
\path{IPRADerivations} polynomial theorem and one bounded trace example through
the complete import--compile--export route.
